\section{Ambientación}

La ambientación del juego se diseñó con fuertes influencias e inspiración en la cultura y el contexto urbano de Ecuador. Se modelaron enemigos basados en agrupaciones delictivas locales y se utilizaron denominaciones características, como se aprecia en la figura 6.

\begin{figure}[H]
    \centering
    % \includegraphics[width=0.6\textwidth]{Imagenes/figura6.png}
    \caption{Tipos de enemigos}
    \label{fig:tipos_enemigos}
\end{figure}

Las zonas urbanas se construyeron utilizando módulos reutilizables de edificaciones con variaciones de color y posición, emulando los sectores populares característicos donde operan estas bandas (figura 7).

\begin{figure}[H]
    \centering
    % \includegraphics[width=0.6\textwidth]{Imagenes/figura7.png}
    \caption{Ambientación de zonas}
    \label{fig:ambientacion_zonas}
\end{figure}

Los elementos estáticos y dinámicos del escenario incluyen vegetación escasa, obstáculos contundentes, cajas y vehículos mal estacionados que el jugador debe esquivar o destruir. El diseño de la vía y los objetos se muestran en las figuras 8 y 9.

\begin{figure}[H]
    \centering
    % \includegraphics[width=0.6\textwidth]{Imagenes/figura8.png}
    \caption{Diseño de la vía}
    \label{fig:ambientacion_calle}
\end{figure}

\begin{figure}[H]
    \centering
    % \includegraphics[width=0.6\textwidth]{Imagenes/figura9.png}
    \caption{Objetos interactivos en el juego}
    \label{fig:objetos_juego}
\end{figure}

Musicalmente, cada zona cuenta con un tema distintivo basado en géneros de rock y sus variantes, aportando identidad y dinamismo a la acción. Asimismo, se integró un sistema de audio reactivo que responde a eventos clave como disparos, destrucción de obstáculos, eliminación de enemigos y pantallas de victoria o derrota.
